\scanpage{64}
% archive sheet number: 63
\deleted{\uline{Question}} \inserted{N.B.} Si $\underline{\mathrm{Pic}}_{X/S}$ est simple sur $S$ aux pts de $\underline{\mathrm{Pic}}^0_{X/S}$, \deleted{\unsure{et si}} [ou ce
qui revient au $=$, le long de \uuline{la section} unité \ldots{}]
\deleted{la réunion des \unsure{composantes de} $S$, en \unsure{présence|propres}} \deleted{Alors} alors
$\underline{\mathrm{Pic}}^0_{X/S}$ est \uline{ouvert} dans $\underline{\mathrm{Pic}}_{X/S}$ [c'est \unsure{pour} tt \unsure{préschém.|préschéma} en
groupes loc. de type fini G sur $S$, car s'il est \uline{simple} sur $S$,
alors \deleted{\unsure{en}} comme chaque fois qu'on peut lui \unsure{imposer}
\unsure{connexion|convenablement} \unsure{par|que} les fibres par une section des \unsure{uns}
\unsure{composantes} \inserted{\uline{fibres géom} \deleted{\unsure{irréductibles}} \uline{loc. intègres}} \deleted{\unsure{irréductibles, les pts}},
$G^0$ est ouvert en $G$].
% margin note, linked by a bent arrow to the bracket
\margin{$\underline{\mathrm{IV}}~7.(6.6.\deleted{\unsure{bis|ter}})$}

Supposons le \uline{propre} sur $S$, donc c'est un schém. abélien
sur $S$. Alors (a) les fibres de $\underline{\mathrm{Pic}}_{X/S}$ sont simples
(b) \deleted{\unsure{$\Omega$}}$\unsure{Q|\mathrm O|\mathcal O}_X$ est coh. plat sur $S$ en \deleted{\unsure{codim}} \uline{dim $1$}
\inserted{i.e. $\underline{\omega}_{\underline{\mathrm{Pic}}_{X/S}}$ est loc. libre \ldots}

\deleted{[ \unsure{Réuni|Réun.|Réc.}.} Est-il vrai réciproquement que (a) et (b)
impliquent que $\underline{\mathrm{Pic}}_{X/S}$ est simple sur $S$ le long de
la section unité ??? Noter qu'en car. résiduelles $0$,
(a) est automatique, et \deleted{\unsure{je}} je pense que (b)
implique en effet (a), dans la forme plus générale :
Si $G$ est un schém en groupes sur $S$ \deleted{\unsure{en car. $0$}}, $S$ à
corps résiduels de car. $0$, et si \deleted{\unsure{$\Omega$}} $\underline\omega_{G/S}$ [faisceau
conormal à la section unité] est \uline{loc. libre}, \uuline{alors} $G$ est
simple sur $S$ le long de la section unité \ldots\ldots

De plus, il y a lieu de se demander si pour $X/S$
projectif et simple, et la car. $0$, les conditions
qu'on vient de dire sont satisfaites automatiquemt. On
en a une démonstration \uline{purement Algébrique}
via $\underline{\mathrm{Pic}}$ dans le cas où $S$ est \uline{réduit}. Dans le
cas $S$ arbitraire, je ne sais rien, \uline{\unsure{au secours !|au concours !|au contraire !}}
